\begin{table}[H]
\centering
\caption{Robustness: Leave-One-Sector-Out and Placebo Outcomes}
\begin{threeparttable}
\begin{tabular}{lc}
\toprule
\multicolumn{2}{l}{\textit{Panel A: Leave-one-sector-out}} \\
\midrule
Specification & Belgium $\times$ Post \\
\midrule
Excl. Agriculture & -0.0249 \\
Excl. Industry & -0.0193 \\
Excl. Construction & -0.0194 \\
Excl. Trade/Transport/Hosp. & -0.0202 \\
Excl. ICT & -0.0216 \\
Excl. Finance/Insurance & -0.0211 \\
Excl. Real Estate & -0.0293 \\
Excl. Prof./Admin Services & -0.0220 \\
Excl. Public Admin/Educ/Health & -0.0222 \\
Excl. Arts/Other Services & -0.0210 \\
\midrule
Baseline (all sectors) & -0.0221 \\
\midrule
\multicolumn{2}{l}{\textit{Panel B: Placebo outcomes}} \\
\midrule
Public sector (NACE O--Q) & -0.0209 \\
 & (0.0291) \\
Real estate (NACE L) & 0.0424 \\
 & (0.0655) \\
\bottomrule
\end{tabular}
\begin{tablenotes}[flushleft]
\small
\item \textit{Notes:} Panel A shows Belgium $\times$ Post coefficient when excluding one NACE A*10 sector at a time. Panel B tests placebo outcomes: public administration (NACE O--Q) and real estate (NACE L) should be unaffected by employer SSC cuts because public employers do not profit-maximize and real estate is capital-intensive. All specifications include country$\times$sector and time fixed effects with country-clustered standard errors.
\end{tablenotes}
\end{threeparttable}
\label{tab:robustness}
\end{table}
